\documentclass[11pt]{article}
\usepackage{amsmath,amssymb,amsthm}
\usepackage{geometry}
\usepackage{booktabs}
\usepackage{graphicx}
\usepackage{hyperref}
\usepackage{enumerate}
\usepackage{listings}
\usepackage{xcolor}
\usepackage{longtable}

\geometry{letterpaper, margin=1in}

\lstset{
  basicstyle=\ttfamily\small,
  keywordstyle=\color{blue},
  commentstyle=\color{gray},
  stringstyle=\color{red},
  breaklines=true,
  columns=fullflexible,
  keepspaces=true,
  language=C++
}

\title{\textbf{Formation Engine Methodology} \\
\Large Section 8c --- Pipe Flow and Thermal Analysis Module}
\author{Formation Engine Development Team}
\date{Version 3.0.1 --- As Implemented in Codebase}

\begin{document}
\maketitle

\tableofcontents
\newpage

% ============================================================================
\section{Introduction and Motivation}
% ============================================================================

Sections~8 and~8b established the thermodynamic framework and heat-gated
reaction control within the atomistic simulation kernel.  This section
extends the thermal infrastructure to \textbf{macroscopic pipe flow and
heat transfer analysis}, providing a deterministic, interactive module
for studying internal forced convection in cylindrical conduits.

The pipe flow and thermal module occupies Layer~5 (macro-scale) of the
VSEPR-SIM layer stack (cf.\ \texttt{include/layer\_stack.hpp}) and is
designed to:

\begin{enumerate}
  \item Solve steady-state internal pipe flow for both laminar and turbulent
        regimes,
  \item Couple flow solutions to convective heat transfer via established
        Nusselt correlations,
  \item Model thermal conduction through pipe walls,
  \item Propagate energy balance through multi-segment series networks,
  \item Support stochastic network generation for uncertainty studies,
  \item Provide parametric sweep capability for design-space exploration,
  \item Export results as Markdown reports and CSV datasets.
\end{enumerate}

All computations use SI units internally.  Unit conversions are applied
\textbf{only at the display boundary}, preserving kernel clarity and
preventing unit confusion in future extensions.

% ============================================================================
\section{Governing Equations}
% ============================================================================

\subsection{Reynolds Number and Flow Regime Classification}

The Reynolds number characterizes the flow regime:

\begin{equation}
  \mathrm{Re} = \frac{\rho \, v \, D}{\mu}
  \label{eq:reynolds}
\end{equation}

where $\rho$ is the fluid density~[kg/m$^3$], $v$ the mean velocity~[m/s],
$D$ the inner pipe diameter~[m], and $\mu$ the dynamic viscosity~[Pa\,s].

The flow is classified as:

\begin{equation}
  \text{Regime} =
  \begin{cases}
    \text{Laminar}    & \mathrm{Re} < 2300 \\
    \text{Turbulent}  & \mathrm{Re} \geq 2300
  \end{cases}
\end{equation}

\noindent
The transition region ($2000 < \mathrm{Re} < 4000$) is handled by the
Churchill correlation, which provides a smooth blend (Section~\ref{sec:friction}).

\subsection{Pressure Drop: Darcy--Weisbach}

The pressure drop across a pipe segment of length~$L$ is:

\begin{equation}
  \Delta P = f \cdot \frac{L}{D} \cdot \frac{\rho \, v^2}{2}
  \label{eq:darcy-weisbach}
\end{equation}

where $f$ is the Darcy friction factor.

\subsection{Friction Factor}
\label{sec:friction}

\subsubsection{Laminar Regime}

For $\mathrm{Re} < 2300$, the analytical Hagen--Poiseuille result is used:

\begin{equation}
  f = \frac{64}{\mathrm{Re}}
  \label{eq:laminar-friction}
\end{equation}

\subsubsection{Turbulent and Transitional Regime: Churchill Approximation}

The implicit Colebrook--White equation:

\begin{equation}
  \frac{1}{\sqrt{f}} = -2.0 \log_{10}
    \!\left(\frac{\varepsilon/D}{3.7} + \frac{2.51}{\mathrm{Re}\sqrt{f}}\right)
  \label{eq:colebrook}
\end{equation}

is solved using the explicit Churchill approximation~\cite{churchill1977},
which avoids iteration and smoothly covers both the transitional and
fully turbulent regimes:

\begin{align}
  A &= \left[2.457 \ln\!\left(\frac{1}{(7/\mathrm{Re})^{0.9}
       + 0.27\,\varepsilon/D}\right)\right]^{16} \\
  B &= \left(\frac{37530}{\mathrm{Re}}\right)^{16} \\
  f &= 8\left[\left(\frac{8}{\mathrm{Re}}\right)^{12}
       + \frac{1}{(A+B)^{3/2}}\right]^{1/12}
  \label{eq:churchill}
\end{align}

\subsection{Hagen--Poiseuille Volume Flow Rate (Laminar Verification)}

For laminar flow, the analytical volume flow rate provides a cross-check:

\begin{equation}
  Q = \frac{\pi \, D^4 \, \Delta P}{128 \, \mu \, L}
  \label{eq:hagen-poiseuille}
\end{equation}

This identity is validated in the test suite (Section~\ref{sec:validation}).

% ============================================================================
\section{Heat Transfer Model}
% ============================================================================

\subsection{Nusselt Number Correlations}

\subsubsection{Laminar: Constant Wall Temperature}

For fully-developed laminar flow with constant wall temperature:

\begin{equation}
  \mathrm{Nu} = 3.66
  \label{eq:nu-laminar}
\end{equation}

\subsubsection{Turbulent: Dittus--Boelter Correlation}

For turbulent flow ($\mathrm{Re} > 2300$, $0.6 \leq \mathrm{Pr} \leq 160$),
the Dittus--Boelter equation for heating:

\begin{equation}
  \mathrm{Nu} = 0.023 \, \mathrm{Re}^{0.8} \, \mathrm{Pr}^{0.4}
  \label{eq:dittus-boelter}
\end{equation}

where the Prandtl number is:

\begin{equation}
  \mathrm{Pr} = \frac{\mu \, c_p}{k_f}
  \label{eq:prandtl}
\end{equation}

\subsection{Convection Coefficient}

The convective heat transfer coefficient is derived from the Nusselt number:

\begin{equation}
  h = \frac{\mathrm{Nu} \cdot k_f}{D}
  \label{eq:h-conv}
\end{equation}

where $k_f$ is the fluid thermal conductivity~[W/(m\,K)].

\subsection{Thermal Resistance Network}

Each pipe segment is modeled as two thermal resistances in series:

\begin{equation}
  R_{\text{conv}} = \frac{1}{h \cdot \pi \, D_i \, L}
  \label{eq:r-conv}
\end{equation}

\begin{equation}
  R_{\text{wall}} = \frac{\ln(r_o / r_i)}{2\pi \, k_w \, L}
  \label{eq:r-wall}
\end{equation}

\begin{equation}
  R_{\text{total}} = R_{\text{conv}} + R_{\text{wall}}
  \label{eq:r-total}
\end{equation}

where $r_i = D_i/2$, $r_o = D_o/2$, and $k_w$ is the wall material
thermal conductivity.

\subsection{Outlet Temperature: Exponential Model}

For constant wall temperature, the fluid outlet temperature is:

\begin{equation}
  T_{\text{out}} = T_w - (T_w - T_{\text{in}}) \, e^{-\mathrm{NTU}}
  \label{eq:t-out}
\end{equation}

where the Number of Transfer Units is:

\begin{equation}
  \mathrm{NTU} = \frac{h \cdot P \cdot L}{\dot{m} \, c_p}
  \label{eq:ntu}
\end{equation}

with $P = \pi D_i$ (wetted perimeter), $\dot{m} = \rho v A$ (mass flow
rate), and $c_p$ the fluid specific heat.

\subsection{Heat Transfer and LMTD}

The total heat exchanged in a segment is:

\begin{equation}
  \dot{Q} = \dot{m} \, c_p \, (T_{\text{out}} - T_{\text{in}})
  \label{eq:energy-balance}
\end{equation}

The log-mean temperature difference provides the effective driving force:

\begin{equation}
  \Delta T_{\text{lm}} = \frac{\Delta T_1 - \Delta T_2}
                              {\ln(\Delta T_1 / \Delta T_2)}
  \label{eq:lmtd}
\end{equation}

where $\Delta T_1 = T_w - T_{\text{in}}$ and
$\Delta T_2 = T_w - T_{\text{out}}$.

\subsection{Thermal Effectiveness}

The network effectiveness is defined as:

\begin{equation}
  \varepsilon = \frac{T_{\text{out}} - T_{\text{in}}}
                     {T_w - T_{\text{in}}}
  \label{eq:effectiveness}
\end{equation}

representing the fraction of maximum achievable temperature change.

% ============================================================================
\section{Network Topology}
% ============================================================================

\subsection{Deterministic Series Network}

The primary mode solves $N$ identical pipe segments in series.  The outlet
temperature of segment~$i$ becomes the inlet temperature of segment~$i+1$:

\begin{equation}
  T_{\text{in}}^{(i+1)} = T_{\text{out}}^{(i)}, \qquad i = 0, \ldots, N-1
\end{equation}

All segments share the same geometry, material, and flow velocity.
Pressure drops and heat transfer rates are accumulated:

\begin{align}
  \Delta P_{\text{total}} &= \sum_{i=0}^{N-1} \Delta P_i \\
  \dot{Q}_{\text{total}}  &= \sum_{i=0}^{N-1} \dot{Q}_i
\end{align}

\subsection{Stochastic Network}

The stochastic mode introduces controlled randomness via a seeded
\texttt{mt19937\_64} generator:

\begin{itemize}
  \item Segment lengths drawn from a log-normal distribution:
        $L_i \sim \text{LogNormal}(\ln L_0, \, 0.3)$
  \item Inner diameters perturbed:
        $D_i = D_0 + \mathcal{N}(0, \, 0.002\,\text{m})$, clamped to $\geq 5$~mm
  \item Wall thickness preserved from the nominal configuration
\end{itemize}

\noindent
\textbf{Determinism guarantee:} Identical seed $\Rightarrow$ bit-identical
output.  Different seeds $\Rightarrow$ statistically different networks.
This is validated in the test suite (Section~\ref{sec:validation}).

% ============================================================================
\section{Parametric Sweeps}
% ============================================================================

The engine provides three built-in parametric sweep functions, each
varying a single design parameter while holding all others at their
configured values:

\begin{center}
\begin{tabular}{llll}
\toprule
Sweep & Parameter & Range & Outputs \\
\midrule
\texttt{sweep\_velocity}         & $v$ [m/s]   & $[v_{\min}, v_{\max}]$ & $\Delta P$, $\dot{Q}$, $T_{\text{out}}$, $\varepsilon$ \\
\texttt{sweep\_diameter}         & $D_i$ [m]   & $[D_{\min}, D_{\max}]$ & $\Delta P$, $\dot{Q}$, $T_{\text{out}}$, $\varepsilon$ \\
\texttt{sweep\_wall\_temperature} & $T_w$ [K]  & $[T_{\min}, T_{\max}]$ & $\Delta P$, $\dot{Q}$, $T_{\text{out}}$, $\varepsilon$ \\
\bottomrule
\end{tabular}
\end{center}

Each sweep returns $n$ points uniformly distributed across the range.
Sweeps can be exported to CSV or visualized in the GUI panel
(Section~\ref{sec:gui}).

% ============================================================================
\section{Material and Fluid Database}
% ============================================================================

\subsection{Fluid Properties}

\begin{center}
\begin{tabular}{lcccc}
\toprule
Fluid & $\rho$ [kg/m$^3$] & $\mu$ [Pa\,s] & $c_p$ [J/(kg\,K)] & $k_f$ [W/(m\,K)] \\
\midrule
Water (20\textdegree C) & 998.2 & $1.002 \times 10^{-3}$ & 4182 & 0.598 \\
Water (80\textdegree C) & 971.8 & $3.545 \times 10^{-4}$ & 4197 & 0.670 \\
Air (20\textdegree C)   & 1.204 & $1.825 \times 10^{-5}$ & 1007 & 0.0257 \\
Ethylene Glycol          & 1113  & $1.61 \times 10^{-2}$  & 2382 & 0.258 \\
Engine Oil (SAE~30)      & 876   & 0.290                  & 1964 & 0.145 \\
\bottomrule
\end{tabular}
\end{center}

\subsection{Pipe Materials}

\begin{center}
\begin{tabular}{lccccc}
\toprule
Material & $k_w$ [W/(m\,K)] & $\varepsilon$ [m] & $\rho_w$ [kg/m$^3$] & $\sigma_y$ [MPa] & $\alpha$ [$10^{-6}$/K] \\
\midrule
Copper    & 385   & $1.5 \times 10^{-6}$ & 8960 & 210 & 16.5 \\
Steel     & 50.2  & $4.5 \times 10^{-5}$ & 7850 & 250 & 12.0 \\
PVC       & 0.16  & $1.5 \times 10^{-6}$ & 1400 & 45  & 70.0 \\
Stainless & 16.3  & $2.0 \times 10^{-6}$ & 8000 & 205 & 17.3 \\
Aluminum  & 237   & $1.5 \times 10^{-6}$ & 2700 & 276 & 23.1 \\
\bottomrule
\end{tabular}
\end{center}

\noindent
Source: Incropera et al.\ (2007), engineering data tables.  Values represent
typical properties; temperature-dependent lookup is a planned future extension.

% ============================================================================
\section{Unit Conversion Architecture}
% ============================================================================

A core design principle of the pipe thermal module is that \textbf{unit
conversions occur only at the display boundary}.  The kernel operates
entirely in SI:

\begin{center}
\begin{tabular}{ll}
\toprule
Quantity & Kernel Unit \\
\midrule
Temperature  & K (Kelvin) \\
Pressure     & Pa (Pascal) \\
Length        & m (metre) \\
Velocity     & m/s \\
Mass flow    & kg/s \\
Heat flux    & W \\
Conductivity & W/(m\,K) \\
\bottomrule
\end{tabular}
\end{center}

\noindent
The \texttt{pipe\_thermal::units} namespace provides conversion functions
for display:

\begin{lstlisting}
namespace pipe_thermal::units {
  double K_to_C(double K);      double C_to_K(double C);
  double K_to_F(double K);      double F_to_K(double F);
  double Pa_to_kPa(double Pa);  double Pa_to_bar(double Pa);
  double Pa_to_psi(double Pa);  double Pa_to_atm(double Pa);
  double m_to_mm(double m);     double m_to_in(double m);
  double m3s_to_Lmin(double);   double m3s_to_gpm(double);
  double W_to_kW(double W);     double W_to_BTUhr(double W);
}
\end{lstlisting}

These functions are \textbf{never} called inside \texttt{solve\_segment()},
\texttt{solve\_series\_network()}, or \texttt{solve\_stochastic\_network()}.
They are only invoked by the TUI, GUI, report generator, and CSV export.

% ============================================================================
\section{Implementation Architecture}
% ============================================================================

\subsection{File Map}

\begin{center}
\begin{tabular}{ll}
\toprule
File & Purpose \\
\midrule
\texttt{include/pipe\_thermal/pipe\_thermal\_engine.hpp} & Kernel: physics, network solver, sweeps, export \\
\texttt{include/pipe\_thermal/pipe\_thermal\_panel.hpp}  & ImGui GUI panel \\
\texttt{apps/pipe\_thermal\_tui.cpp}                     & Interactive TUI (REPL) \\
\texttt{apps/pipe\_thermal\_viewer.cpp}                  & Standalone ImGui viewer \\
\texttt{tests/test\_pipe\_thermal.cpp}                   & Validation suite (32 assertions) \\
\bottomrule
\end{tabular}
\end{center}

\subsection{Data Flow}

The data flow follows a strict unidirectional pipeline:

\begin{enumerate}
  \item \textbf{Input:} User provides \texttt{PipeThermalConfig} (all SI).
  \item \textbf{Kernel:} \texttt{solve\_segment()} computes all intermediates in SI.
  \item \textbf{Propagation:} \texttt{solve\_series\_network()} chains segments,
        accumulating totals.
  \item \textbf{Collection:} \texttt{NetworkResult} stores all segment results
        and network-level metrics.
  \item \textbf{Display:} TUI/GUI/report calls \texttt{units::*} for human-readable output.
\end{enumerate}

\noindent
This ensures that no intermediate computation ever encounters a unit ambiguity.

\subsection{Anti-Black-Box Design}

Every intermediate quantity is stored in \texttt{SegmentResult} and fully
accessible to the user:

\begin{itemize}
  \item Geometry: $L$, $D_i$, $D_o$, wall thickness, flow area, wetted perimeter
  \item Flow: $v$, $\mathrm{Re}$, laminar/turbulent flag, $f$, $\Delta P$, $Q$, $\dot{m}$
  \item Thermal: $\mathrm{Nu}$, $h$, $R_{\text{conv}}$, $R_{\text{wall}}$, $R_{\text{total}}$,
        $T_{\text{in}}$, $T_{\text{out}}$, $T_w$, LMTD, $\dot{Q}$
  \item Derived: $\varepsilon$, pumping power
\end{itemize}

No hidden state.  No implicit assumptions.

% ============================================================================
\section{TUI Application}
% ============================================================================

The TUI (\texttt{apps/pipe\_thermal\_tui.cpp}) provides a REPL interface:

\begin{lstlisting}[language=bash]
$ pipe-thermal-tui
> c              # show configuration
> f 2            # select fluid (Air)
> m 3            # select material (PVC)
> v 2.5          # set velocity to 2.5 m/s
> d 0.050        # set inner diameter to 50 mm
> r              # run series network
> s              # run stochastic network
> sv 0.5 5.0 20  # sweep velocity from 0.5 to 5.0 m/s, 20 points
> report         # generate Markdown report
> csv            # generate CSV export
> q              # quit
\end{lstlisting}

Batch mode is also supported:

\begin{lstlisting}[language=bash]
$ pipe-thermal-tui --batch 10
\end{lstlisting}

This solves a 10-segment series network with default parameters and
prints a complete Markdown report to stdout.

% ============================================================================
\section{GUI Panel}
\label{sec:gui}
% ============================================================================

The GUI panel (\texttt{include/pipe\_thermal/pipe\_thermal\_panel.hpp})
provides an ImGui-based interactive interface with:

\begin{itemize}
  \item Sliders for all geometric and flow parameters
  \item Combo boxes for fluid and pipe material selection
  \item Auto-solve on parameter change
  \item Per-segment detail table with all intermediates
  \item Parametric sweep visualization via \texttt{ImGui::PlotLines}
\end{itemize}

The panel is conditionally compiled under \texttt{BUILD\_VISUALIZATION}
and can be embedded in any ImGui host application.

% ============================================================================
\section{Validation}
\label{sec:validation}
% ============================================================================

The test suite (\texttt{tests/test\_pipe\_thermal.cpp}) contains 32~assertions
across 7~test groups:

\subsection{Test Group 1: Laminar Hagen--Poiseuille}

\begin{itemize}
  \item $\mathrm{Re} < 2300$ confirmed
  \item $f = 64/\mathrm{Re}$ verified to $\pm 1\%$
  \item $\Delta P$ matches analytical Hagen--Poiseuille within $1\%$
  \item $\mathrm{Nu} = 3.66$ verified to $\pm 0.01$
\end{itemize}

\subsection{Test Group 2: Turbulent Dittus--Boelter}

\begin{itemize}
  \item $\mathrm{Re} > 2300$ confirmed
  \item $\mathrm{Nu} = 0.023 \, \mathrm{Re}^{0.8} \, \mathrm{Pr}^{0.4}$ verified to $\pm 5\%$
  \item $h = \mathrm{Nu} \cdot k_f / D$ verified to $\pm 1\%$
\end{itemize}

\subsection{Test Group 3: Energy Balance}

The energy balance $\dot{Q} = \dot{m} \, c_p \, \Delta T$ is verified to
$\pm 0.1\%$ for both laminar and turbulent cases.

\subsection{Test Group 4: Unit Conversions}

All 14~conversion functions are verified against exact reference values.

\subsection{Test Group 5: Deterministic Reproducibility}

\begin{itemize}
  \item Same seed $\Rightarrow$ identical $T_{\text{out}}$ to machine precision
  \item Different seed $\Rightarrow$ different $T_{\text{out}}$
\end{itemize}

\subsection{Test Group 6: Sweep Consistency}

Pressure drop increases monotonically with velocity across the sweep range.

\subsection{Test Group 7: Welford Heat Capacity}

A synthetic Gaussian energy distribution ($\mu = 1000$, $\sigma = 50$,
$100{,}000$ samples) is fed to both the naive and Welford estimators.
The Welford result is verified to match the naive result within $10\%$
and to agree with the analytical $C_v = \sigma^2 / (k_B T^2)$.

\subsection{Validation Summary}

All 32~assertions pass.  The validation covers the complete physics chain
from Reynolds number through outlet temperature to network effectiveness.

% ============================================================================
\section{Prerequisite Fixes}
% ============================================================================

Three prerequisite issues were resolved as part of this development phase:

\subsection{Coulomb Stability (lj\_coulomb.cpp)}

The Coulomb potential $U = k_e q_i q_j / r$ has a $1/r$ singularity that
causes integrator instability at short range.  A softcore damping term was
introduced:

\begin{align}
  r_{\text{soft}}^2 &= r^2 + \delta^2 \qquad (\delta = 0.2\,\text{\AA}) \\
  U_{\text{coul}} &= \frac{k_e \, q_i \, q_j}{\sqrt{r_{\text{soft}}^2}} \\
  F_{\text{coul}} &= \frac{k_e \, q_i \, q_j \, r}{r_{\text{soft}}^2 \cdot \sqrt{r_{\text{soft}}^2}}
\end{align}

Additionally, a force magnitude clamp ($F_{\max} = 1000$~kcal/(mol\,\AA))
prevents catastrophic acceleration if particles overlap.

\subsection{Neighbor List (neighbor\_list.hpp)}

The O($N^2$) placeholder was replaced with a cell-list accelerated Verlet
neighbor list:

\begin{itemize}
  \item Cell grid with side length $\geq r_{\text{list}} = r_{\text{cut}} + r_{\text{skin}}$
  \item 13-cell half-shell stencil (avoids double counting)
  \item Periodic boundary condition support via minimum image convention
  \item Half-skin rebuild criterion: rebuild when any particle has moved
        more than $r_{\text{skin}}/2$ since last build
\end{itemize}

Build cost: O($N$) for uniform systems.

\subsection{Welford Heat Capacity (thermodynamics.hpp)}

The \texttt{heat\_capacity\_welford()} function was added to the
\texttt{atomistic::thermo} namespace, using the existing
\texttt{atomistic::OnlineStats} (Welford algorithm) for numerically
stable online variance computation:

\begin{equation}
  C_v = \frac{\mathrm{Var}(E)}{k_B \, T^2}
  \label{eq:welford-cv}
\end{equation}

This avoids the catastrophic cancellation that can occur with the
naive two-pass formula when $\langle E^2 \rangle \approx \langle E \rangle^2$.

% ============================================================================
\section{References}
% ============================================================================

\begin{enumerate}
  \item Incropera, F.P., DeWitt, D.P., Bergman, T.L., and Lavine, A.S.
        (2007). \textit{Fundamentals of Heat and Mass Transfer}. 6th~ed.
        Wiley.

  \item White, F.M. (2011). \textit{Fluid Mechanics}. 7th~ed. McGraw-Hill.

  \item Colebrook, C.F. (1939). ``Turbulent Flow in Pipes, with Particular
        Reference to the Transition Region Between the Smooth and Rough
        Pipe Laws.'' \textit{J.\ Inst.\ Civil Eng.}, \textbf{11}, 133--156.

  \item Churchill, S.W. (1977). ``Friction-Factor Equation Spans All
        Fluid-Flow Regimes.'' \textit{Chem.\ Eng.}, \textbf{84}(24), 91--92.
        \label{churchill1977}

  \item Dittus, F.W.\ and Boelter, L.M.K. (1930). ``Heat Transfer in
        Automobile Radiators of the Tubular Type.''
        \textit{Univ.\ Calif.\ Publ.\ Eng.}, \textbf{2}(13), 443--461.

  \item Welford, B.P. (1962). ``Note on a Method for Calculating
        Corrected Sums of Squares and Products.''
        \textit{Technometrics}, \textbf{4}(3), 419--420.
\end{enumerate}

\end{document}
